\documentclass[12pt,a4paper]{article}

\usepackage[a4paper,margin=1in]{geometry}
\usepackage[T1]{fontenc}
\usepackage[utf8]{inputenc}
\usepackage{enumitem}
\setlength{\emergencystretch}{3em}

\title{\textbf{Supervisor Meeting Minutes}\\Meeting 3}
\author{Present: Matthew Chalmers, Jonas Leipner}
\date{Date: 13 October 2026}

\begin{document}

\maketitle

\section{Achieved in Previous Week}

\noindent\textit{Previous week: 6--13 October 2026}
\begin{itemize}[leftmargin=*]
    \item Re-evaluated the constraints affecting dynamic code loading in production.
    \item Investigated the distinction between over-the-air patching and true runtime code loading.
\end{itemize}

\section{Issues Discussed \& Solutions Agreed}

\begin{itemize}[leftmargin=*]
    \item \textbf{Shorebird limitation:} Shorebird can deliver code patches over the air, but the application must restart. It does not load new code dynamically inside an active game session.
    \item \textbf{True dynamic loading:} Runtime scripting through Lua or JavaScript, or the \texttt{dart\_eval} Dart bytecode interpreter, is required for code to execute dynamically during a session.
    \item \textbf{\texttt{dart\_eval} architecture:} \texttt{dart\_eval} is sandboxed. An internal bridge is required to pass Flame game state into the sandbox, calculate values such as combat damage, and return validated results to Flame.
    \item \textbf{Research scope:} The project should compare plugin architectures with runtime dynamic code loading approaches used in commercial games.
\end{itemize}

\section{Objectives for This Week}

\noindent\textit{Upcoming week: 13--20 October 2026}
\begin{itemize}[leftmargin=*]
    \item Test communication between the Flame Engine and the \texttt{dart\_eval} sandbox.
    \item Research memory optimisation and dynamic loading practices in commercial game development.
\end{itemize}

\section*{\small AI-use Disclaimer}
\small The meeting content and source notes for these minutes are genuine. Artificial intelligence was used only to organise and summarise the notes into a more digestible format; it was not used to invent the meeting content.

\end{document}
